\chaptersubtitle{przez wielce uczonego Rogera Bacona spisana, \\ mała książeczka o Alchemii.}

\chapter{Zwierciadło Alchemii}

\section{O definicjach Alchemii}

W starożytnych księgach znaleźć można wiele definicji owej sztuki, których zamysł winniśmy rozważyć w tym rozdziale. 

Hermes\footnote{\textbf{Hermes Trismegistos} - Hermes Po-Trzykroć-Wielki. Bóstwo hellenistyczne, któremu przypisuje się autorstwo części traktatów alchemicznych, w tym \textit{Corpus Hermeticum} (zob. \textit{Corpus Hermeticum} tłum. R. Wąsowski, 2019), oraz \textit{Tablic Szmaragdowych} (zob. \textit{Hermetyzm} R. Bugaj, Wrocław, 1991, s. 120). [przyp. tłum.]} mówi o tej nauce: \alchquote{Alchemia jest nauką o ciałach, w prosty sposób z jednego i przez jedno złożoną, łączy ona rzeczy cenniejsze poprzez poznanie i działanie oraz przekształca je naturalnym zmieszaniem w doskonalszy rodzaj}.  Ktoś inny zaś powiada: \alchquote{Alchemia jest nauką, która uczy, jak przekształcać metal dowolnego rodzaju w dowolny inny}, a to przy pomocy odpowiedniego medykamentu, jak wynika z licznych ksiąg filozofów.

Zatem Alchemia jest nauką, która uczy tworzyć oraz mieszać pewien medykament zwany \textit{Eliksirem}, który, rzucony na metale lub niedoskonałe ciała, czyni je doskonałymi w momencie rzutu.


\section{O zasadach naturalnych i zradzaniu minerałów}

Po drugie, dokładnie przedstawię zasady naturalne i powstawanie minerałów:  po pierwsze należy zauważyć, że minerałami pierwszymi\footnote{\textit{principia mineralia} (łac.), \textit{natural principles} (ang.) -- dosł. zasady mineralne, podstawy mineralne. Chodzi o wcześniej wspomniane \textit{zasady naturalne}. [przyp. tłum.] } w kopalniach są: \textit{Argentum vivum}\footnote{\textit{Argentum vivum} (łac.), \textit{Argent vive} (ang.) -- dosł. żywe srebro, rtęć. [przyp. tłum.]} oraz siarka. Z nich to zrodzone są wszystkie metale i minerały, których wiele jest różnych gatunków. Lecz powiadam, natura zawsze pragnie i dąży do doskonałości złota, lecz różne są akcydensy zmieniające metale, co wyraźnie widać w księgach filozofów. Zgodnie bowiem z czystością i nieczystością dwóch wspomnianych wcześniej, to jest \textit{Argentum vivum} i siarki, powstają czyste i nieczyste metale: złoto, srebro, cyna, ołów, miedź i żelazo. O ich naturze, to jest o czystości i nieczystości, czyli nadwyżce brudu lub niedomiarze, przyjmij za prawdę to, co następuje:

\subsection{O naturze złota}
Złoto jest ciałem doskonałym. Z \textit{Argentum vivum}, bez skazy\footnote{\textit{p\=urum} (łac.), \textit{pure} (ang.) - aby odróżnić tę czystość i czystość \textit{mundus} (łac.), \textit{clean} (ang.), zdecydowano się pierwsze z nich tłumaczyć na: \textit{bez skazy}, drugie zaś na: \textit{czysty}. [przyp. tłum.]}, stałego, jasnego\footnote{\textit{claro} (łac), \textit{clear} (ang.) - jasny, przejrzysty, klarowny. [przyp. tłum.]} oraz czerwonego, i z siarki czystej, stałej, czerwonej i niepalnej, pochodzi, oraz skaz żadnych nie ma.

\subsection{O naturze srebra}
Srebro jest ciałem czystym, bez skazy i prawie doskonałym. Z \textit{Argentum vivum}, bez skazy, prawie stałego, jasnego i białego oraz z podobnej siarki pochodzi. Niczego, poza odrobiną stałości, koloru i ciężkości, mu nie potrzeba.

\subsection{O naturze cyny}
Cyna jest ciałem czystym i niedoskonałym. Z \textit{Argentum vivum} bez skazy, stałego i niestałego, jasnego, zewnętrznie białego, lecz wewnątrz czerwonego oraz z podobnej siarki pochodzi. Potrzebuje jedynie wywarzenia lub strawienia.

\subsection{O naturze ołowiu}
Ołów jest ciałem brudnym i niedoskonałym. Z \textit{Argentum vivum} nieczystego, niestałego, ziemistego, mętnego, zewnętrznie nieco białego, lecz wewnątrz czerwonego i siarki po części palnej pochodzi. Pragnie usunięcia skaz, stałości, koloru i wypalenia.

\subsection{O naturze miedzi}
Miedź jest ciałem brudnym i niedoskonałym. Z \textit{Argentum vivum} skalanego, niestałego, ziemistego, palnego, czerwonego i mętnego oraz z podobnej siarki pochodzi. Pragnie usunięcia skaz, utrwalenia i ciężkości oraz ma zbyt wiele nieczystego  koloru i ziemistości niepalnej.

\subsection{O naturze żelaza}
Żelazo jest ciałem nieczystym i niedoskonałym. Z \textit{Argentum vivum} skalanego, zbyt stałego, ziemisto-palnego, białego i czerwonego, mętnego  oraz z podobnej siarki pochodzi. Pragnie zjednania, oczyszczenia i ciężkości, a także ma zbyt wiele stałej i nieczystej siarki oraz ziemistości palnej. 
\begin{alchline}
To, co zostało powiedziane, każdy alchemik musi pilnie przestrzegać.
\end{alchline}

\section{O rzeczach z których substancja Eliksiru winna być wydobyta}

Zradzanie metali, zarówno tych doskonałych, jak i niedoskonałych, zostało wystarczająco wytłumaczone przez to, co zostało już powiedziane. Powróćmy zatem do niedoskonałej substancji, którą musimy wybrać i uczynić doskonałą. 

Mając na uwadze, że w poprzednich rozdziałach nauczeni zostaliśmy, że wszelkie metale zrodzone są z \textit{Argentum vivum} i siarki, oraz że ich nieczystości i skazy je gorszą, i że nic nie może być zmieszane z metalami, co by w \textit{Argentum vivum} i siarce pochodzenia swego nie miało. Pozostaje więc wystarczająco jasne, że żadna obca rzecz, która nie ma początku w nich, nie jest w stanie sama ani ich udoskonalić, ani zmienić, ani przekształcić w nowe. 

Zastanawiającym byłoby więc, gdyby jakiś mędrzec skierował swe myśli ku żywym stworzeniom lub roślinom, które są od nich wielce odległe, podczas gdy minerały są im bliskie. Nie sposób też myśleć, że któryś z filozofów umieściłby sztukę we wspomnianych odległych rzeczach, chyba że w celu porównania. 

W każdym razie, ze wspomnianych wcześniej dwóch wszystkie metale pochodzą, i ani nic się do nich nie przylega, ani się z nimi nie łączy, ani ich nie przekształca, prócz tego, co z nich pochodzi. Dlatego też słusznie musimy wziąć \textit{Argentum vivum} i siarkę za substancję naszego kamienia, lecz ani samo \textit{Argentum vivum}, ani siarka sama nie zradzają metali, lecz dopiero z ich połączenia zradzane są różnorodne metale i minerały. Zatem i nasza substancja winna być dobrana z połączenia tych dwóch.

A oto nasza ostateczna tajemnica jest najznamienitsza i najskrytsza: z jakiego minerału, bliższego niż inne, substancja naszego kamienia powinna być wydobyta, a w tym wyborze musimy być bardzo ostrożni. 

Stawiam zatem sprawę jasno, że jeśli substancja ta miałaby być wydobyta z roślin (do których należą: zioła, drzewa oraz wszystko, co z ziemi wyrasta), to musielibyśmy najpierw wydobyć \textit{Argentum vivum} i siarkę przez długie warzenie, z czego zostaliśmy zwolnieni, gdyż natura obdarowała nas \textit{Argentum vivum} i siarką. 

Jeśli zaś mielibyśmy wydobyć ją z istot żywych (przez co mam na myśli: krew, włosy, mocz, odchody, jaja kurze i wszystko inne, co od żywych istot pochodzi), to podobnie przez warzenie musielibyśmy wydobyć \textit{Argentum vivum} i siarkę, z czego również jesteśmy zwolnieni, tak jak byliśmy w poprzednim przypadku.

Moglibyśmy też wydobyć ją z minerałów pośrednich (do których należą: wszelkie magnesje i markasyty, tutia, miedź, ałuny, baurach, sole i wiele innych), z których, jak w poprzednich przypadkach, musielibyśmy wydobyć \textit{Argentum vivum} i siarkę przez warzenie, z czego, podobnie jak w poprzednich przypadkach, zostaliśmy zwolnieni. 

%najprawdopodobniej chodzi o czystą siarkę (zob. G. Starkey, \textit{The True Oil of Sulphur}, w \textit{Collectanea Chemica}, A.E. Waite, 1893).

Gdybyśmy wzięli jednego z siedmiu duchów, to jest: \textit{Argentum vivum} samo, albo siarkę samą, albo \textit{Argentum vivum} i jedną z dwóch siarek, albo \textit{Sulphur vivum\footnote{\textit{Sulphur vivum} (łac.), \textit{Sulphur vive} (ang.) - dosł. żywa siarka. [przyp. tłum.]}}, aurypigment, \textit{Arsenicum citrinum}\footnote{\textit{Arsenicum citrinum} (łac.), \textit{Citrine Arsenicum} (ang) - arszenik cytrynowy. [przyp. tłum]}, albo czerwień samą\footnote{Najpewniej mowa o \textit{realgarze}, który występuje w rudach z aurypigmentem i charakteryzuje się czerwonym kolorem. [przyp. tłum]}, albo im podobne, to niemożliwym byłoby zmienić go, ponieważ natura nigdy nie czyni doskonałym bez równego zmieszania obu, a więc i my nie możemy. Z tych to zatem, jak i z \textit{Argentum vivum} oraz siarki w ich naturze, jesteśmy zwolnieni. Bo gdybyśmy zdecydowali się wybrać je, to musielibyśmy zmieszać je w odpowiednich częściach, których ludzki rozum nie zna, a następnie warzyć je, aż do ścięcia w stałą masę. Ponieważ zaś nie znamy tych części, a znaleźć możemy ciała, w których znajdujemy wspomniane rzeczy w odpowiednich częściach, ścięte i zebrane razem, jesteśmy zwolnieni z konieczności brania obu w ich własnej naturze.

Zachowajcie tę tajemnicę bardziej tajemną: złoto jest doskonałym ciałem męskim, bez żadnych nadmiarów i ubytków, i gdyby przez samo stopienie doskonaliło niedoskonałe ciała, byłoby \textit{Eliksirem} do czerwieni. Srebro natomiast jest prawie doskonałym ciałem żeńskim, i gdyby przez samo stopienie niemalże doskonaliło niedoskonałe ciała, byłoby \textit{Eliksirem} do bieli, czym nie jest i być nie może, ponieważ tylko to, co doskonałe, może być doskonałe. A jeśli doskonałość ta mogłaby być zmieszana z niedoskonałością, to niedoskonałe nie powinno być doskonalone przez doskonałe, lecz raczej ich doskonałość powinna zostać zmniejszona przez niedoskonałość i sama winna stać się niedoskonała. Lecz gdyby były one więcej niż doskonałe, czy to w podwójnej, poczwórnej, stukrotnej, czy w większej proporcji, mogłyby wówczas należycie udoskonalić to, co niedoskonałe. A ponieważ natura zawsze działa w prosty sposób, doskonałość, która w złocie i srebrze jest, jest prosta, niepodzielna i niemieszalna, więc nie można ich też sztuką włożyć w kamień jako ferment dla skrócenia dzieła i w ten sposób przywrócić je do ich pierwotnego stanu, ponieważ to, co najbardziej lotne, pokonuje to, co najbardziej stałe. A że złoto jest ciałem doskonałym, składającym się z \textit{Argentum vivum}, czerwonego i jasnego oraz podobnej siarki, dlatego nie wybieramy go na substancję naszego kamienia do czerwonego \textit{Eliksiru}, ponieważ jest ono po prostu doskonałe, i to bez sztucznego oczyszczania, oraz tak silnie strawione i nakarmione naturalnym gorącem, że naszym sztucznym ogniem ledwo jesteśmy w stanie je lub srebro obrabiać. I choć natura doskonali rzeczy, to jednak nie potrafi ich całkowicie oczyścić, to znaczy: udoskonalić i pozbawić skaz, ponieważ pracuje tylko na tym, co ma. Jeśli zatem wybralibyśmy złoto lub srebro na substancję naszego kamienia, z trudem znaleźlibyśmy choćby nikły ogień działający w nich. I chociaż nie jesteśmy nieświadomi tego ognia, to niestety nie potrafimy ani złota, ani srebra w pełni oczyścić i udoskonalić z powodu ich ścisłego zespolenia i naturalnego składu. Zatem jesteśmy zwolnieni z brania zbyt czerwonego pierwszego i zbyt białego drugiego, gdyż znaleźć możemy rzecz lub ciało z podobnie czystymi, a nawet czystszymi \textit{Argentum vivum} i siarką, na które natura prawie nic nie zdziałała, a które przy pomocy naszego sztucznego ognia i doświadczenia w naszej sztuce jesteśmy w stanie doprowadzić do właściwego strawienia, oczyszczenia, koloru i utrwalenia w nieustannej zręcznej pracy nad nimi. 

Musi zatem zostać wybrana taka substancja, w której znajduje się \textit{Argentum vivum}, czyste, bez skazy, przejrzyste, białe i czerwone, nie w pełni doskonałe, lecz równo i proporcjonalnie zmieszane należytym sposobem z podobną siarką, ścięta w stałą masę, abyśmy dzięki naszej mądrości i rozwadze oraz naszemu sztucznemu ogniowi mogli osiągnąć jej najwyższą czystość i nieskazitelność i doprowadzić ją do takiego stanu, że po zakończeniu pracy będzie tysiąc tysięcy razy silniejsza i doskonalsza niż proste substancje strawione przez ich naturalne ciepło. 

Bądźcie więc mądrzy: jeśli bowiem będziecie subtelni i bystrzy w Rozdziałach moich, w których to przejrzystą prozą jasno wyłożyłem sprawę kamienia łatwą do poznania, zakosztujecie tej zachwycającej rzeczy, w której zawarta jest cała intencja filozofów.

\section{O zasadach pracy z ogniem, jego podtrzymywaniu i kontrolowaniu}

Wierzę, że odkryłeś, jeśli nie jesteś otępialcem, ignorantem lub głupcem, dzięki dotychczas wypowiedzianym słowom, właściwą substancję błogosławionego kamienia uczonych filozofów, na której działa Alchemia, kiedy to niedoskonałe staramy się uczynić doskonałym, i to za pomocą rzeczy bardziej niż doskonałej. I skoro natura dała nam niedoskonałe tylko z doskonałym, naszym to celem jest uczynić substancję, w poprzednich rozdziałach nam przedstawioną, bardziej niż doskonałą przez naszą sztuczną pracę. 

A jeśli nie znamy sposobu działania, to czemu nie patrzymy w jaki sposób natura, która od dawna udoskonala metale, nieustannie działa? Czyż nie widzimy, że w kopalniach, przez nieustanny żar, który w górach panuje, gęsta woda tak się warzy i gęstnieje, że z czasem staje się \textit{Argentum vivum}? I że z tłuszczu ziemi przez  ten sam gorąc i warzenie zradza się siarka! I że poprzez ten sam skwar, nieustannie w nich trwający, zradzają się z nich wszystkie metale, zależnie od ich czystości i nieczystości? I że natura przez same warzenie czyni wszystkie metale, i te doskonałe, i niedoskonałe?

O jakże wielkie szaleństwo! Powiedzcież proszę, co zmusza was do próby doskonalenia rzeczy wcześniej wspomnianych przez melancholijne i fantastyczne sposoby? Mówią przecie: \alchquote{Biada wam, którzy chcecie przewyższyć naturę i uczynić metale doskonalszymi nowymi sposobami lub pracą zrodzoną z waszych bezmyślnych głów. Bóg dał naturze prostą drogę, to jest ciągłe warzenie, a wy, głupcy, gardzicie nią lub nie znacie jeszcze jej}. I nauczają: \alchquote{Ogień i azot wam wystarczą}, \alchquote{Ciepło doskonali wszystko}, oraz: \alchquote{Gotuj, gotuj, gotuj i nie przemęczaj się}, i: \alchquote{Niech ogień będzie delikatny i łagodny oraz zawsze równy, by mógł stale płonąć, a nie pozwól, by zwiększył się, bo jeśli zwiększy się, poniesiesz olbrzymie straty}. I również: \alchquote{Wiedz, że w jednej rzeczy, to jest w kamieniu, jednym sposobem, to jest warzeniem, i w jednym naczyniu całe magisterium się zawiera}. I w innym miejscu mówią: \alchquote{Cierpliwie i ciągle}, oraz: \alchquote{Skrusz go siedmiokrotnie}, i: \alchquote{Kruszone jest ogniem}. I gdzie indziej: \alchquote{Ta praca podobna jest do stworzenia człowieka}, bo podobnie jak niemowlę karmione jest lekkimi daniami, a gdy kości jego się wzmocnią, zaczyna przyjmować coraz to cięższe pokarmy. Podobnie magisterium najpierw potrzebuje małego ognia, z którym należy pracować w każdym procesie warzenia. I chociaż zawsze mówimy o łagodnym ogniu, to jednak naprawdę sądzimy, że w pracy ogień musi być stopniowo i kolejno 

\section{O jakości naczynia i pieca}

Zasady i sposoby pracy już ustaliliśmy, teraz powinniśmy pomówić o naczyniu i piecu: w jaki sposób i z czego powinny być zrobione. Natura bowiem przez naturalny ogień warzy metale w kopalniach, a więc nie dopuszcza, by podobne warzenie odbywało się bez odpowiedniego do tego naczynia. A jeśli chcielibyśmy naśladować naturę w przygotowaniach, to dlaczego mielibyśmy odrzucać jej naczynie? Zobaczmy zatem, gdzie zradzanie metali ma miejsce. 

Po miejscach występowania minerałów wyraźnie widać, że u podnóża góry panuje nieustannie równomierne ciepło, którego naturą jest ciągłe dążenie ku górze, a w dążeniu tym zawsze osusza i sprawia, że krzepnie gęstsza lub grubsza woda, ukryta w łonie oraz żyłach ziemi albo góry, w \textit{Argentum vivum}. A jeśli mineralna tłustość tego miejsca, ogrzana w ten sposób, zbierze się w żyłach ziemi, przepływa ona przez górę i staje się siarką. I można zobaczyć, że we wspomnianych wcześniej żyłach siarka zrodzona z tłustości ziemi spotyka się z \textit{Argentum vivum} i zradza gęstość wody mineralnej. Tam też, przez ciągłe i równomierne ciepło w długim procesie, rodzą się różne metale, zgodnie z różnorodnością miejsca. I w tych miejscach mineralnych znajdziesz ciągłe ciepło! Należy więc odnotować, że zewnętrzna góra mineralna jest wszędzie zamknięta i kamienista, bo gdyby ciepło mogło wydostać się, nigdy nie powstałby żaden metal.

Jeśli zatem zamierzamy naśladować naturę, powinniśmy mieć piec podobny do gór, lecz nie w wielkości, a w trwałości ciepła, tak aby umieszczony w nim ogień, wznosząc się, nie znajdował ujścia, ale by ciepło mogło uderzać w naczynie szczelnie zamknięte, zawierające w sobie substancję kamienia. Naczynie to musi być okrągłe, z małą szyjką, wykonane ze szkła lub z jakiejś ziemi posiadającej jego naturę, to jest zwięzłość. Otwór naczynia należy zamknąć podobnym przykryciem i uszczelnić albo zapieczętować bitumem. I jak w kopalniach, ciepło nie dotyka bezpośrednio \textit{Argentum vivum} i siarki, ponieważ ziemia góry znajduje się wszędzie pomiędzy. Podobnie i ogień nie może bezpośrednio dotykać naczynia zawierającego wspomniane wcześniej rzeczy, ale musi zostać ono umieszczone w innym naczyniu, szczelnie zamkniętym w podobny sposób, aby umiarkowane ciepło mogło dotykać substancji powyżej i poniżej oraz wszędzie, gdzie się ona znajduje, w sposób bardziej odpowiedni i właściwy. Z tego też powodu Arystoteles w \textit{Lumen Luminum}\footnote{zob. \textit{Lumen Luminum} w \textit{Syntagma harmoniae chymico-philosophicae}, Frankfurt 1625, str. 93-106 [przyp. tłum.]} pisze: \alchquote{Merkuriusz przygotowany winien być w potrójnym naczyniu z twardego szkła lub (co lepsze) z ziemi posiadającej naturę szkła}.

\section{O przypadkowych i istotnych kolorach w dziele występujących}

Skoro więc istota kamienia została przedstawiona, powinieneś poznać właściwy sposób pracy, dzięki któremu kamień podczas warzenia mieni się różnymi kolorami. O nich to mówią: \alchquote{Ile kolorów, tyle i nazw}. Zgodnie z tą wielością, nazwy używane przez filozofów są różne: \alchquote{W pierwszej fazie naszego kamienia, fazie zwanej rozkładem, nasz kamień czernieje}. O czerni tej mówią: \alchquote{Kiedy zastaniesz go czarnego, wiedz, że w tej czerni biel jest skryta. A jeśliby ta czerń była najsubtelniejsza z subtelnych, wtedy musisz ją z tej najsubtelniejszej czerni wydobyć}. 

Po \textit{rozkładzie} kamień czerwienieje, lecz nie jest to prawdziwa czerwień. O niej to mówią: \alchquote{Często jest czerwony, często też cytrynowy, często się topi i często krzepnie przed prawdziwą bielą, a następnie rozpuszcza się i ścina, rozkłada się i barwi, uśmierca się i ożywia, czerni się i bieli, czerwieni się, ale również zieleni}. O zieleni zaś mówią: \alchquote{Mieszaj go, dopóki nie będzie zdawał się zielony, wtedy jest duszą}, czy też: \alchquote{Wiedz, że w tej zieleni dusza ta ma panowanie}.

Przed bielą ukazuje się również pawi kolor. Powiadają o nim: \alchquote{Wszystkie kolory świata, które są do pomyślenia, pojawiają się przed bielą, a po nich prawdziwa biel przychodzi}, o której to piszą: \alchquote{Kiedy  zostanie wywarzony  bezbłędnie i czysto, tak że będzie lśnił niczym oko rybie, wtedy spodziewaj się jego użyteczności}. Do tego momentu kamień powinien skrzepnąć okrągle.

Dalej powiadają: \alchquote{Kiedy znajdziesz biel u szczytu naczynia, bądź pewien, że w tej bieli czerwień jest skryta}, i tę czerwień musisz wydobyć. Zatem warz ją, aż cała stanie się czerwona. Lecz pomiędzy prawdziwą bielą a prawdziwą czerwienią znajduje się popielaty kolor. O nim powiadają: \alchquote{Po bieli nie możesz zbłądzić, bo gdy zwiększysz ogień, popielaty kolor Ci się ukaże}, oraz: \alchquote{Nie gardź popiołem, gdyż Bóg winien zwrócić Ci go płynnym, a wtem ostatni Król zostanie z woli Bożej ukoronowany czerwoną koroną}.

\section{O rzutowaniu medykamentu na dowolne niedoskonałe ciało}

W większości dotrzymałem mojej obietnicy znamienitego mistrzostwa w tworzeniu najwspanialszego z \textit{Eliksirów}, czerwonego i białego. Na koniec zajmiemy się rzutowaniem, które jest zakończeniem dzieła, upragnioną i wyczekiwaną radością. 

Czerwony \textit{Eliksir} cytryni się nieskończenie i przemienia każdy metal w najczystsze złoto. Biały \textit{Eliksir} nieskończenie bieli i czyni każdy metal doskonale białym. Ale wiemy, że jedne z metali są dalsze od doskonałości niż inne, i że jedne są bliższe niż inne. I mimo że każdy metal może być przez \textit{Eliksir} doprowadzony do doskonałości, to najbliższe są doprowadzane łatwiej, szybciej i doskonalej niż te, które są dalsze. A kiedy natrafiamy na metal, który jest prawie doskonały,  nie musimy zajmować sobie głowy tymi, które są daleko. 


Jeśli chodzi o to, które metale są bliskie, a które dalekie, oraz który z nich jest najbliższy doskonałości, jeśli jesteś rozumny, powinieneś dostrzec to jasno i prawdziwie wyłożone w poprzednich rozdziałach. Bez wątpliwości, ten, kto spoglądając w moje \textit{Zwierciadło}, wykazał się wielką bystrością, przez swą pracowitość znajdzie prawdziwą substancję i w pełni pozna, na które z ciał medykament winien być rzucony, by uczynić je doskonałym. 

Pierwsi tej sztuki, którzy odkryli ją przez swoją filozofię, wskazują palcem prostą drogę, kiedy mówią: \alchquote{Natura zawiera naturę. Natura przezwycięża naturę. I Natura spotykając swoją naturę, niezmiernie się raduje i przemienia w inne natury}, oraz w innym miejscu: \alchquote{Każdy raduje się w podobnych sobie, bo podobieństwo jest przyczyną przyjaźni}. 

Z nich też wielu pozostawiło znamienne tajemnice: \alchquote{Wiedz, że dusza szybko wchodzi w swoje ciało, które nie sposób połączyć z innym ciałem}, a w innym miejscu: \alchquote{Bo dusza szybko wchodzi w swoje ciało, a gdy połączy się z innym ciałem utracisz całą pracę, gdyż bliskość sama w sobie jest jaśniejsza}.

I ponieważ materialne rzeczy w tym sposobie pracy stają się niematerialne, a niematerialne stają się materialne, u kresu dzieła całe ciało staje się duchowo stałe. I ponieważ duchowy \textit{Eliksir}, niezależnie czy biały, czy czerwony, jest tak wielce przygotowany i wywarzony poza swoją własną naturę, że nie jest dziwnym, iż nie może być zmieszany z ciałem, na które jest rzutowany, przez samo tylko stopienie. Trudem również jest rzutować go na tysiąc tysięcy i więcej, by natychmiast przeniknąć i przekształcić je.

Zdradzę Ci teraz tajemnicę wspaniałą i szczególnie skrytą: jedna część musi zostać zmieszana z tysiącem części najbliższego ciała, a następnie szczelnie zamknięta w odpowiednim naczyniu i ustawiona na piecu, najpierw na małym ogniu, który powinien być zwiększany w przeciągu trzech dni, dopóki zawartość nie stanie się nierozerwalnie złączona. I to jest praca trzech dni. Po tym czasie każda część musi być rzucona na kolejne tysiące części najbliższego ciała. A to jest praca na kolejny dzień, bądź godzinę, albo i moment, za co nasz wspaniały Bóg wiecznie winien być sławiony!
